\chapter{Stage III：exact WKB と大域接続}

\solproblem{III.1　Riccati 係数}{
対数 WKB ansatz を代入し最初の4係数と枝の parity を求めよ。}{
$[-\hbar^2\partial_x^2+Q(x)]\psi=0$ に
$\psi=\rme^{\hbar^{-1}\int^xS(x',\hbar)\dd x'}$ を代入すると
\begin{equation}
 S^2+\hbar S'=Q,
 \qquad S=\sum_{n\geq0}\hbar^nS_n.
\end{equation}
従って $S_0^2=Q$ と
\begin{equation}
 S_n=-\frac{1}{2S_0}
 \left(S_{n-1}'+\sum_{j=1}^{n-1}S_jS_{n-j}\right)
\end{equation}
を得る。$S_0^\pm=\pm\sqrt Q$ とすれば
\begin{align}
 S_1&=-\frac{Q'}{4Q},\\
 S_2^\pm&=\pm\left(
 \frac{Q''}{8Q^{3/2}}-\frac{5(Q')^2}{32Q^{5/2}}\right),\\
 S_3&=-\frac{Q'''}{16Q^2}
 +\frac{9Q'Q''}{32Q^3}-\frac{15(Q')^3}{64Q^4}.
\end{align}
sheet 交換で $S_{2j}$ は符号を変え、$S_{2j+1}$ は不変。
$\vsub{S}{odd}=(S^+-S^-)/2$ が偶数次数、
$\vsub{S}{even}=(S^++S^-)/2$ が奇数次数を含み、
$\vsub{S}{even}=-(\hbar/2)\partial_x\log\vsub{S}{odd}$ である。}

\solproblem{III.2　最も単純な lateral Borel 和}{
$\sum_{n\geq0}n!\hbar^{n+1}$ の Borel 変換、二つの lateral 和、差を求めよ。}{
Borel 変換は
\begin{equation}
 \widehat F(t)=\sum_{n\geq0}\frac{n!}{n!}t^n=\frac{1}{1-t}.
\end{equation}
正実軸上の $t=1$ に pole があるので
\begin{equation}
 \mathcal S_{0\pm}F(\hbar)
 =\int_{0}^{\infty\rme^{\pm\rmi0}}
 \frac{\rme^{-t/\hbar}}{1-t}\dd t
\end{equation}
と処方する。実部は principal value、虚部は半留数である。上側 minus
下側を本問の規約とすれば
\begin{equation}
 \mathcal S_{0+}F-\mathcal S_{0-}F
 =2\pi\rmi\rme^{-1/\hbar}.
\end{equation}
符号は lateral path の命名を逆にすれば反転するが、ambiguity の大きさと
action $A=1$ は不変である。}

\solproblem{III.3　単純 turning point と Airy 接続}{
単純 turning point を Airy 方程式へ写し、局所接続行列を求めよ。}{
$Q(x_0)=0$、$Q'(x_0)\neq0$ とし
\begin{equation}
 \zeta(x)=\left[\frac32\int_{x_0}^x\sqrt{Q(t)}\dd t\right]^{2/3},
 \qquad \psi=(\zeta/Q)^{1/4}w
\end{equation}
と置く。leading order では
$\hbar^2w_{\zeta\zeta}=\zeta w$、解は
$\operatorname{Ai}(\zeta/\hbar^{2/3})$ と $\operatorname{Bi}$ である。
dominant/subdominant の係数列を $(\vsub{a}{dom},\vsub{a}{sub})^{\mathsf T}$ と
し、Stokes 線を正向きに横切る規約では dominant 係数は不変で
subdominant 係数が $\pm\rmi$ 倍だけ跳ぶ：
\begin{equation}
 \binom{\vsub{a}{dom}}{\vsub{a}{sub}}_{\!\mathrm{after}}
 =\begin{pmatrix}1&0\\ \pm\rmi&1\end{pmatrix}
 \binom{\vsub{a}{dom}}{\vsub{a}{sub}}_{\!\mathrm{before}}.
\end{equation}
符号は $\sqrt Q$ の sheet と横断方向で決まる。ordering を
(sub,dom) に変えれば転置位置へ移る。いずれも determinant 1 で
Wronskian を保存する。}

\solproblem{III.4　Stokes 行列の path ordering}{
三つの triangular Stokes 行列と propagation 行列を掛け、非可換性を示せ。}{
\begin{equation}
 U(a)=\begin{pmatrix}1&a\\0&1\end{pmatrix},\quad
 L(b)=\begin{pmatrix}1&0\\b&1\end{pmatrix},\quad
 P(p)=\begin{pmatrix}p&0\\0&p^{-1}\end{pmatrix}
\end{equation}
とする。各 determinant は1なので、任意の path-ordered product も
determinant 1 で Wronskian を保存する。しかし
\begin{equation}
 U(a)L(b)=\begin{pmatrix}1+ab&a\\b&1\end{pmatrix},\qquad
 L(b)U(a)=\begin{pmatrix}1&a\\b&1+ab\end{pmatrix},
\end{equation}
また $PUP^{-1}=U(p^2a)$ である。従って、例えば
$U(a_3)L(a_2)P(p)U(a_1)$ は横断順序を入れ替えた積と異なる。
determinant だけでは順序誤りを検出できないため、各 crossing の sheet、
向き、basis ordering を ledger に記録する必要がある。}

\solproblem{III.5　barrier resonance と Gamow 幅}{
二 turning point barrier の leading outgoing 条件と幅を導け。}{
classically allowed well の一周期 action を $I(E)$、escape barrier の
Euclidean action を
$K(E)=\int_{x_1}^{x_2}\sqrt{2m(V-E)}\dd x>0$ とする。Airy 接続を一周
掛けると、leading real quantization は
\begin{equation}
 I(E_n)=2\pi\hbar(n+1/2).
\end{equation}
barrier の小さい transmission probability は
$T(E)=\rme^{-2K(E)/\hbar}$。classical period
$\vsub{T}{cl}=\partial I/\partial E$ ごとに一回 escape を試みるなら
decay rate は $T(E_n)/\vsub{T}{cl}$、従って
\begin{equation}
 \Gamma_n=\frac{\hbar}{\vsub{T}{cl}(E_n)}
 \rme^{-2K(E_n)/\hbar}
\end{equation}
（左右二出口なら各寄与を加える）。all-order exact WKB では $K$ を
Borel resummed open Voros period、$I$ を closed quantum period に置き換え、
Stokes multiplier も含む $\Cincoming(E)=0$ を解く。}

\solproblem{III.6　inverted Rosen--Morse の pole}{
hypergeometric 接続から incoming 係数と共鳴・反共鳴枝を求めよ。}{
$V=U_0/\cosh^2(\beta x)$、$E=\hbar^2k^2/(2m)$、
$s(s+1)=-2mU_0/(\beta^2\hbar^2)$ とする。$\xi=\tanh\beta x$ により
associated Legendre 方程式となり、$x\to+\infty$ で outgoing な解を
$x\to-\infty$ へ接続すると、非零位相を除き
\begin{equation}
 \vsub{A}{in}(k)=
 \frac{\Gamma(1-\rmi k/\beta)\Gamma(-\rmi k/\beta)}
 {\Gamma(-\rmi k/\beta-s)
  \Gamma(-\rmi k/\beta+s+1)}.
\end{equation}
Gamma 関数は零を持たないので、分母の pole が $\vsub{A}{in}=0$ を与える。
$\kappa=[8mU_0/(\beta^2\hbar^2)-1]^{1/2}$ とすれば resonance branch は
\begin{equation}
 \vsup{k_n}{R}=\frac{\beta}{2}[\kappa-
 \rmi(2n+1)],\quad
 \vsup{E_n}{R}=\frac{\hbar^2\beta^2}{8m}
 [\kappa-
 \rmi(2n+1)]^2.
\end{equation}
real potential の反共鳴は
$\vsup{k_n}{AR}=-(\vsup{k_n}{R})^*$。どちらも同じ解析式の別 sheet である。}

\solproblem{III.7　単純零点と二重零点の Laurent 展開}{
$S(E)=N(E)/\Cincoming(E)$ を単純零点と二重零点の周りで展開せよ。}{
$x=E-E_*$ とする。単純零点
$\Cincoming=c_1x+c_2x^2/2+\cdots$、
$N=N_0+N_1x+\cdots$ なら
\begin{equation}
 S(E)=\frac{N_0}{c_1}\frac1x
 +\left(\frac{N_1}{c_1}-\frac{N_0c_2}{2c_1^2}\right)+\order(x).
\end{equation}
二重零点
$\Cincoming=c_2x^2/2+c_3x^3/6+\cdots$ なら
\begin{equation}
 S(E)=\frac{2N_0}{c_2}\frac1{x^2}
 +\left(\frac{2N_1}{c_2}-\frac{2N_0c_3}{3c_2^2}\right)\frac1x
 +\order(1).
\end{equation}
Stage VII の Jordan resolvent の $x^{-2}$ 項と $x^{-1}$ 項がこの二係数に
対応する。ただし行列問題では determinant だけでなく左右 null vector
を含む residue を計算する。}

\solproblem{III.8　cubic metastable potential の Stokes graph}{
二枚の $\sqrt Q$ sheet 上で Stokes graph と outgoing path を構成せよ。}{
例 $Q(x,E)=x^2-gx^3-2E$ の三つの turning point $x_j$ を求め、各点から
\begin{equation}
 \im\int_{x_j}^{x}\sqrt{Q(t,E)}\dd t=0
\end{equation}
を満たす三本の Stokes 曲線を数値積分する。branch cut を turning point
間に置き、sheet $+$ では $S_0=+\sqrt Q$、cut 横断後は sheet $-$ と記録する。
origin 側の regular/subdominant 解から metastable well を通り、barrier
turning points の Airy 行列、cut 交換、外側の Stokes jump を順に掛け、
外側で incoming 係数を0にする path を選ぶ。正解の検査は
(i) 各 crossing の順序が連続、(ii) branch 交換が偶数回なら元 sheet、
(iii) 総行列 determinant が1、(iv) contour の小変形で $\Cincoming$ が不変、
の四点である。図の見た目だけで sheet を判定してはならない。}

\solproblem{III.9　analytic coordinate rotation と Voros period}{
turning points、cycle、積分路を同時に回転したとき quantum period が不変であることを示せ。}{
$x=\rme^{\rmi\theta}y$ と同時に potential parameter と turning points を
解析的に移し、cycle $\gamma$ を $\gamma_\theta$ へ写す。Riccati 1-form は
座標変換則により
\begin{equation}
 S_x(x,E;\hbar)\dd x=S_y(y,E;\hbar)\dd y+dd F
\end{equation}
となる。closed cycle では exact differential の積分は0なので
\begin{equation}
 \VorosPeriod_\gamma=oint_\gamma\vsub{S}{odd}\dd x
 =\oint_{\gamma_\theta}\vsub{S}{odd,\theta}\dd y.
\end{equation}
open path でも endpoint 規格化を同時に変換すれば同じである。変形領域を
pole、turning point、branch cut、Borel singular direction が横切ると
homotopy class または lateral sum が変わり、留数あるいは Stokes
automorphism を追加しなければならない。}

\solproblem{III.10　piecewise constant barrier の transfer matrix}{
transfer-matrix pole と WKB 接続条件を比較せよ。}{
区間 $j$ の波を $A_j\rme^{\rmi k_jx}+B_j\rme^{-\rmi k_jx}$ と書く。
幅 $d_j$ の propagation と境界 matching は
\begin{equation}
 P_j=\begin{pmatrix}\rme^{\rmi k_jd_j}&0\\0&\rme^{-\rmi k_jd_j}\end{pmatrix},
\quad
 I_{j,j+1}=\frac12
 \begin{pmatrix}1+k_j/k_{j+1}&1-k_j/k_{j+1}\\
                 1-k_j/k_{j+1}&1+k_j/k_{j+1}\end{pmatrix}.
\end{equation}
$M=I_{N,N+1}P_N\cdots P_1I_{0,1}$ を path order で作る。左右とも
outgoing の条件を課すと、採用した係数順序に応じて $M_{22}(E)=0$
（または等価な minor）が pole 条件である。電流規格化した interface
行列では $\det M=1$ となり、$|\det M-1|$ を丸め誤差 monitor にする。
smooth barrier の区間数を増やすと、$\log P_j$ の和は WKB propagation
action へ、interface 行列は turning-point connection へ収束する。
pole の比較は同じ $k_j$ branch と outgoing convention で行う。}
